\origpage{326}{N. Robertson, P.D. Seymour / Journal of Combinatorial Theory, Series B 92 (2004) 325--357}
It was shown in [3] that

\textbf{1.2.} \textit{For every graph $H$, if $G$ has no minor isomorphic to $H$, then every ``highly connected component'' of $G$ can ``almost'' be drawn on a surface on which $H$ cannot be drawn.}

(The meanings of ``highly connected component'' and ``almost'' here are complicated and we shall postpone the exact statement of this theorem as long as possible. Surfaces are connected and compact.)

We may assume the surface in 1.2 is without boundary; and since up to homeomorphism there are only finitely many such surfaces in which $H$ cannot be drawn, to prove 1.1 and hence Wagner's conjecture it suffices to show that

\textbf{1.3.} \textit{If $\Sigma_1, \ldots, \Sigma_n$ are surfaces then for every infinite set $\mathcal{F}$ of graphs, if every highly connected component of every member of $\mathcal{F}$ can almost be drawn in one of $\Sigma_1, \ldots, \Sigma_n$, then some member of $\mathcal{F}$ is isomorphic to a minor of another member of $\mathcal{F}$.}

To prove 1.3 we use the main results of two other papers of this series [4,5]. The main result of [4] asserts that, if $\mathcal{F}$ is an infinite set of graphs and all the highly connected components of all members of $\mathcal{F}$ have a certain ``well-behaved'' structure, then some member of $\mathcal{F}$ is isomorphic to a minor of another member of $\mathcal{F}$. It therefore suffices to show that the hypothesis of 1.3 implies that all these highly connected components have a well-behaved structure. To show this, we apply the main result of [5], which asserts that for any infinite set of hypergraphs all drawable in a fixed surface (where the edges of the hypergraphs all have two or three ends, and each edge is labeled from a fixed well-quasi-order), some member of the set is isomorphic to a minor of another (with an appropriate definition of ``minor'' for hypergraphs).

In Sections 2--10 we finish the proof of Wagner's conjecture, and in Section 11 we prove a slight strengthening.

\sectionhead{2. Hypergraphs and tangles}

For the purposes of this paper, a hypergraph $G$ consists of a finite set $V(G)$ of \textit{vertices}, a finite set $E(G)$ of \textit{edges}, and an incidence relation between them. The vertices incident with an edge are the \textit{ends} of the edge (A hypergraph is thus a graph if every edge has one or two ends.) A hypergraph $H$ is a \textit{subhypergraph} of a hypergraph $G$ (written $H \subseteq G$) if $V(H) \subseteq V(G)$, $E(H) \subseteq E(G)$, and for every $v \in V(G)$ and $e \in E(H)$, $e$ is incident with $v$ in $G$ if and only if $v \in V(H)$ and $e$ is incident with $v$ in $H$. If $G_1, G_2$ are subhypergraphs of $G$ we denote by $G_1 \cup G_2, G_1 \cap G_2$ the subhypergraphs with vertex sets $V(G_1) \cup V(G_2)$, $V(G_1) \cap V(G_2)$ and edge sets $E(G_1) \cup E(G_2)$, $E(G_1) \cap E(G_2)$, respectively. A \textit{separation} of $G$ is an ordered pair $(G_1, G_2)$ of subhypergraphs with $G_1 \cup G_2 = G$ and $E(G_1 \cap G_2) = \emptyset$, and its \textit{order} is $|V(G_1 \cap G_2)|$.

A central idea in our approach is that of a \textit{tangle} in a hypergraph, which was introduced in [2]. Intuitively, a tangle of order $\theta$ is a ``$\theta$-connected component'' of the hypergraph, which therefore resides on one side or the other of every separation of order $< \theta$. Formally, let $G$
